\documentclass[10pt,a4paper]{article}
\pdfinfoomitdate=1
\pdftrailerid{}
\pdfsuppressptexinfo=-1
\usepackage[margin=29mm]{geometry}
\usepackage[T1]{fontenc}
\usepackage{lmodern,microtype,amsmath,amssymb,amsthm,booktabs,array,graphicx}
\pdfmapfile{+lm.map}
\pdfmapfile{+symbols.map}
\PassOptionsToPackage{hyphens}{url}
\usepackage[hidelinks]{hyperref}
\hypersetup{
 pdftitle={Faster Algorithms for Subset Sum using Filtered Representation Trees},
 pdfauthor={},pdfsubject={},pdfkeywords={},pdfcreator={},pdfproducer={}
}
\usepackage[numbers,sort&compress]{natbib}
\setlength{\emergencystretch}{1.5em}
\Urlmuskip=0mu plus 1mu
\newtheorem{theorem}{Theorem}[section]
\newtheorem{lemma}[theorem]{Lemma}
\theoremstyle{definition}
\newtheorem{assumption}[theorem]{Assumption}
\newcommand{\HH}{\mathsf H}
\newcommand{\pos}[1]{[#1]_+}
\newcommand{\poly}{\operatorname{poly}}
\newcommand{\supp}{\operatorname{supp}}
\newcommand{\E}{\mathbb E}
\newcommand{\Prob}{\mathbb P}
\title{\bfseries Faster Algorithms for Subset Sum \\using Filtered Representation Trees}
\author{Samuel J. Adams\textsuperscript{1} \quad Luis Iglesias\textsuperscript{1}\\[0.35em]
{\small\textsuperscript{1}PQProtect Inc., USA}\\[0.35em]
{\footnotesize\texttt{samuel.j.adams80@gmail.com}\quad\texttt{luisiglesias.vigo@outlook.com}}}
\date{\today}
\begin{document}
\maketitle

\begin{abstract}
We study filtered representation algorithms for random planted subset sum
on $n$ input numbers at density one. Within the binary representation-tree
frameworks compared here, we conditionally improve the classical running
time from $2^{0.2690n+o(n)}$ of
Howgrave-Graham~\cite[Theorem~7]{HG26} to $2^{0.2634n+o(n)}$.
With quantum access to classical memory (QRACM), the running time improves
from $2^{0.2356n+o(n)}$~\cite[Section~4.3 and Table~2]{BBSS} to
$2^{0.2100n+o(n)}$. With quantum access to quantum memory (QRAQM),
it improves from $2^{0.2156n+o(n)}$~\cite[Section~5.3]{BBSS} to
$2^{0.1980n+o(n)}$. For each new construction, total space satisfies
the same upper bound as running time. The constructions combine
coordinate-dependent representation laws with type-class filtering and
cost models adapted to their search engines. The quantum-memory
construction uses one common size for all nonroot lists, so its remaining
resource rates follow from short entropy formulas. We derive every
source, index, candidate and reporting cost and verify the three finite
profiles using directed interval calculations. The classical and
quantum-search bounds require assumptions on recursive list
distributions, survival and access. The quantum-memory bound further
requires a coherent-update and hard-workspace interface, whose validity
is not implied by an ordinary expected update count.
\end{abstract}

\let\thefootnote\relax\footnotetext{Supplementary material is available at \texttt{https://github.com/pqprotect/subset-sum}.}

\section{Introduction}
Subset sum asks which numbers in a given collection add up to a specified
target. The question is elementary, but large instances can be difficult:
the number of possible subsets grows exponentially with the input size.
The problem is a basic example in the study of exact algorithms, and its
random instances also provide a useful test for techniques developed in
cryptanalysis. Improvements often carry over to related search problems,
including decoding and learning problems.

A central idea is to give a solution many descriptions. Instead of
splitting a binary solution into two disjoint parts, a representation
algorithm allows partial vectors with additional digits, including
negative ones. Many pairs of partial vectors can then describe the same
solution. A tree of such representations creates room for extra modular
constraints, which keep the intermediate lists manageable. The challenge
is to preserve enough descriptions of a solution while paying for every
partial vector and every pair that is tested.

Filtering addresses a second source of work. Two partial vectors may
satisfy a modular equation but still have incompatible coordinate
patterns. A filter organizes the lists so that fewer such pairs are
examined. Its value depends on how it is used: building an index,
reporting an entire output list and finding a single partner have
different costs. Quantum search changes some of these costs, while
leaving the cost of writing a classical table unchanged. This paper
studies that interaction for three computational models.

\subsection{Related work}
Representation algorithms for hard random subset sum were introduced by
Howgrave-Graham and Joux~\cite{HGJ}, who obtained running time
$2^{0.3370n+o(n)}$
by allowing many decompositions of a solution. Becker, Coron and
Joux~\cite{BCJ} introduced signed representations and obtained numerical
time and space bounds of approximately $2^{0.2909n+o(n)}$. Their Theorem~2 also
gives a formal almost-all-instances guarantee: for every fixed
$\varepsilon>0$, a fraction at least $1-140\cdot2^{-\varepsilon n}$ of
their equibalanced instances of density below one admits the modified
algorithm in time $\widetilde O(2^{(0.2910+3\varepsilon)n})$.
That guarantee concerns their instance distribution and construction;
it does not establish the class-dependent load assumptions used here.

Under a standard list heuristic, Bonnetain, Bricout, Schrottenloher and
Shen (BBSS)~\cite{BBSS} extended
the representation alphabet to include the digit $2$ and obtained
classical running time $2^{0.2830n+o(n)}$. Their quantum algorithms distinguish two
memory models. In QRACM, a quantum computation can read a classical
stored table coherently; their quantum-search construction gives time
and space $2^{0.2356n+o(n)}$~\cite[Section~4.3 and Table~2]{BBSS}.
In QRAQM, the stored data may also be quantum. Their representation walk
runs in time $2^{0.2156n+o(n)}$ and uses space $2^{0.2110n+o(n)}$
under an additional
update heuristic~\cite[Section~5.3]{BBSS}. Their alternative construction
removes that heuristic with time and space $2^{0.2182n+o(n)}$, while
retaining the classical list heuristic~\cite[Section~6.4]{BBSS}.
BBSS explicitly distinguish a correct list count from the stronger
requirements of an efficient coherent walk update.

The directly related classical filtering result is Howgrave-Graham's
analysis of type-class bucketing in representation trees~\cite{HG26}.
Its heuristic Theorem~7 gives running time $2^{0.2690n+o(n)}$ at depth
five. The filters
avoid generating many modularly matched pairs with the wrong split
pattern; their benefit makes deeper trees useful. Our classical result
extends this filtered binary-tree construction by choosing the useful
and filter laws jointly across two coordinate classes. The use of
representations, type-class filters and deeper trees is already present
in this earlier work.

Directed indices and implicit queries also occur in the lattice
algorithms of Hhan, Kim, Kim and Laarhoven~\cite[Lemmas~8.3 and~8.4]{HKKL}.
They motivate the distinction between data writes, query search and
all-output reporting used here. Their lattice distribution and success
arguments do not establish the corresponding properties of subset-sum
lists. Our precise task is planted balanced modular subset sum; the comparison
concerns the random density-one representation frameworks and keeps
their respective guarantees distinct. Other analyses report smaller
running times, including quantum time $2^{0.2090n+o(n)}$
in~\cite{LiLi}, and classical
time $2^{0.2400n+o(n)}$ and quantum time $2^{0.1843n+o(n)}$
in~\cite{Li}. Their representation-counting and ternary-join
analyses do not supply the particular binary-tree cost and access
interfaces used here. We do not claim an unrestricted record
across those different analyses.

\subsection{Contributions}
We give three profiled constructions and complete conditional analyses
for random planted balanced modular subset sum at density one. The
precise problem distribution is defined in Section~\ref{sec:preliminaries}.
Time counts preprocessing and table construction; space counts all
exponentially large lists, indices and temporary storage. The assumptions
are stated formally before the cost derivations.

\begin{theorem}[Classical bound]\label{thm:classical}
Under the recursive-list and survival assumption
(Assumption~\ref{ass:lists}) and the classical access conditions in
Assumption~\ref{ass:access}, the specified classical construction has
constant success probability and time and total space
\[
 T_C,S_C\le 2^{0.2634n+o(n)}.
\]
\end{theorem}

This extends the usual random-list heuristic to the actual
class-dependent, filtered recursive lists: modular loads, candidate
counts and solution survival must hold jointly for that construction.
The premise has the same purpose as the classical list hypotheses in
BBSS and Howgrave-Graham, but we do not prove that their formulations
imply ours. The independent uniform input coefficients, fresh split
offsets and uniformly sampled source subsets provide concrete sources
of randomization. The checked capacities exclude profiles that demand
more vectors than the underlying type domains contain. Exact type
counts and an explicit covering construction
justify the local probabilities. Once the actual candidate counts are
given, classical indexing, bucket scans and deduplication are constructive.
The unresolved part is the distribution and survival of the correlated
recursive lists. The theorem therefore remains conditional.

\begin{theorem}[QRACM bound]\label{thm:qracm}
Under Assumption~\ref{ass:lists} and the reversible padded-query,
occurrence and reporting conditions in Assumption~\ref{ass:access},
the specified QRACM construction has constant success probability and
\[
 T_Q,S_Q<2^{0.2100n+o(n)}.
\]
All exponentially large persistent storage is classical.
\end{theorem}

This construction requires no coherent updates of its stored lists.
Its additional issue is preparing and searching the correct implicit
occurrence states. The required padding bounds and occurrence control
are properties of the actual filtered queries, and are not consequences
of amplitude amplification alone. The product-type structure and exact
static bucket navigation provide concrete parts of the access
implementation. Uniform padding and recursive survival remain
unproved profile-specific requirements. We do not claim that the new
QRACM premises follow merely from the list heuristic used in the BBSS
construction.

\begin{theorem}[QRAQM bound]\label{thm:qraqm}
Under Assumption~\ref{ass:lists}, the coherent canonical setup and access
conditions in Assumption~\ref{ass:access}, and the effective-update and
hard-workspace conditions in Assumption~\ref{ass:updates}, the specified
QRAQM construction has constant success probability and
\[
 T_W,S_W<2^{0.1980n+o(n)}.
\]
\end{theorem}

The quantum-memory statement uses a substantially stronger hypothesis
than static list counting. A source exchange must admit a canonical
coherent update, its inverse, the required root test and extraction,
with controlled error and a hard space bound. The charged update
exponent is suggested by typed expected-work estimates. Turning those
estimates into the full operational interface is an outstanding step;
an ordinary classical mean does not prove it. The assumption is
analogous in purpose to the update heuristic used for the BBSS running time
$2^{0.2156n+o(n)}$, but no equivalence is established. It is an additional
requirement relative to their construction running in time
$2^{0.2182n+o(n)}$, which avoids that heuristic. Thus the faster running time relies on a stronger unresolved quantum
implementation requirement.

The three proofs appear in Sections~\ref{sec:classical},
\ref{sec:qracm} and~\ref{sec:qraqm}. Each gives the algorithm, the full
resource expressions and a finite rational profile with positive margins.
Exact type-flow checks and directed interval calculations establish the
numerical inequalities; they do not
validate the distributional or coherent-update assumptions. No empirical
validation of those assumptions is claimed.

\subsection{Techniques}
\paragraph{Coupling different coordinate laws in one search.}
A single coordinate distribution need not serve the whole vector. We use
two equally sized coordinate groups and allow them to favor different
digits and split patterns. Both groups still belong to the same partial
vector and the same modular equation. Their combined cost can therefore
balance expensive and inexpensive representation choices before the
largest cost is selected. In the two root subtrees the profiles are
exchanged, preserving a shared whole-vector description. This enlarges the profile family available to the optimization. The
reported bounds use jointly chosen useful splits and filters in that
family. They change the probability laws of the construction, not only
the number of modular bits at each level.

\paragraph{Pricing data writes and query search separately.}
A useful-pair predicate can be symmetric even when an efficient quantum
filter should be directed. Copying an entry into more data buckets costs
linear work. Visiting more query addresses can instead benefit from
quantum search. We choose the two incidence conditions with these
separate prices in mind: accepting a larger query search can be worthwhile
if it reduces the number of stored copies. This read/write distinction
is inherited from directed filtering techniques such as~\cite{HKKL}.
Here it is integrated into every relevant subset-sum merge and into the
joint choice of representation laws. The index writes remain charged
in full.

\paragraph{Choosing which lists to materialize.}
The three search engines pay for different objects. The classical
algorithm constructs both incidence tables and checks every candidate.
The QRACM construction materializes sibling subtrees once and searches
one implicit path using nested amplitude amplification. Its stored merges
must still report all their outputs. In QRAQM, a changed source element
triggers searches against maintained data indices. Separating these
operations prevents a one-solution search cost from being used for the
construction of an entire list, or an expected output count from being
used as the cost of finding the outputs. The resulting accounting permits
different useful filters and rate balances in the three models.

\paragraph{Preserving global types through small local filters.}
Small product filters make incidence decoding efficient, but requiring
one fixed pattern in every block can discard valid global
representations. We use a subexponential collection of compatible block
patterns and verified local codebooks. Every useful global pair has a
pattern in that collection, and its canonical first common address can
be tested efficiently. The product-code principle is already present
in~\cite[Proposition~5]{HG26}; the construction here gives simultaneous
coverage of every prescribed global useful pair. This makes the static
filter implementation explicit without treating recursive list
independence as a consequence of local coverage.

\paragraph{A common list height for the quantum walk.}
Every nonroot list in the selected QRAQM profile has the same size
exponent. The modular constraints are chosen to preserve that size from
one merge to the next. The propagation accounting then has no ratio between
parent and child list sizes, and the parameter choice becomes easy to
describe: the height and update rate come from two
information quantities at one merge level, and all other rates follow.
This is a property of the selected witness, not a general theorem that
equal heights are always optimal. Nor does it control update tails or
supply a coherent implementation; those are precisely the extra
requirements in Theorem~\ref{thm:qraqm}.

The constants are certified at finite probability laws. Exact zero
channel cells are retained as zeros, type-flow identities are checked
over the rationals, and nonstructural inequalities have fixed positive
margins. This makes the claimed bounds independent of optimizer
trajectories and displayed decimal rounding, while keeping numerical
feasibility distinct from the assumptions needed by the algorithms.

\subsection{Outlook}
The three constructions show how a more flexible representation profile
and careful model-dependent filtering costs can improve the exponents
within this framework. The principal open questions concern the
assumptions: proving load and survival estimates for the filtered
recursive lists, and constructing the canonical coherent updates needed
for the quantum-memory target. The explicit profiles provide fixed
objects for such analysis and for experimental study. Further numerical
optimization may refine the constants, but resolving these proof
obligations would give a more substantial improvement in the strength
of the results.



\section{Preliminaries}\label{sec:preliminaries}
\subsection{Problem, rates and notation}
For even $n$, put $M=2^n$. Choose $a_1,\ldots,a_n$ independently and
uniformly from $\mathbb Z_M$, choose an independent uniform binary vector
$e$ of weight $n/2$, and put $t=\langle a,e\rangle\pmod M$.
Given $(a,t)$, the task is to return any binary solution. A rate $c$
means $2^{cn+o(n)}$ operations or stored words of polynomial bit length.
Depths, alphabets and numbers of coordinate classes are fixed as $n$ grows.
Classical memory permits ordinary random access. QRACM permits coherent
reads of classical tables; QRAQM also permits coherent quantum stored data.

All logarithms are base two, $\HH$ is Shannon entropy, and $0\log0=0$.
For a fixed coordinate partition, $W$ denotes the class of a uniform
coordinate. Conditional entropies given $W$ are weighted averages over
those classes. A prescribed empirical distribution is called a type.


\subsection{Counting and quantum algorithm primitives}
\label{sec:imported-primitives}

We record the counting and quantum algorithm results used in the cost
analysis. A quantum cost includes each indicated operation and its inverse;
an addressable table is charged separately when it is constructed or stored.
These results do not assert that the recursively generated subset-sum lists
have the load or coherence properties required by those operations.

\begin{lemma}[Type-class sizes {\cite[Theorem~11.1.3]{CT06}}]
\label{lem:type-size}
Let $p$ be an $n$-type on a fixed finite alphabet $\mathcal A$.
For its type class $T_n(p)$,
\[
 (n+1)^{-|\mathcal A|}2^{n\HH(p)}
 \leq |T_n(p)|\leq 2^{n\HH(p)}.
\]
In particular, $|T_n(p)|=2^{n\HH(p)+O(\log n)}$.
For a fixed number of coordinate classes, the product of their type classes
has rate $\HH(X\mid W)$, where $W$ identifies the class.
\end{lemma}
The conditional pair counts used below follow by applying the multinomial
formula within each class and, when the parent is fixed, within each parent
symbol. We give these particular counts in Lemma~\ref{lem:type-counts};
they do not require a distributional assumption on algorithmically
generated lists.

\begin{lemma}[Amplitude amplification
  {\cite[Theorems~3 and~4]{BHMT}}]
\label{lem:amplitude-amplification}
Let a measurement-free algorithm $A$ prepare a state whose projection onto
a reversibly testable good subspace has squared norm $p>0$.
When $p$ is known, amplitude amplification uses $O(p^{-1/2})$ applications
of $A$, $A^{-1}$ and the good-state test to prepare its normalized good
component coherently, using the exact final rotation of Theorem~4.
When $p$ is unknown, the measured search of Theorem~3 finds a good outcome
using an expected $O(p^{-1/2})$ such applications.
\end{lemma}
The coherent statement requires the preparation circuit, its inverse and
the good-state phase test. It preserves relative amplitudes of accepted
addresses. Thus uniform candidate addresses give uniform accepted
\emph{occurrences}; different occurrences can still encode the same vector.
The measured unknown-$p$ algorithm is not a clean state-preparation
interface for a nested quantum procedure. The access assumption below
specifies the preparations actually needed in that setting.

\begin{lemma}[Quantum reporting on an implicit tree
  {\cite[Lemma~8.4]{HKKL}}]
\label{lem:tree-reporting}
Let a fixed binary tree have $S$ vertices, depth $2^{o(n)}$ and
$D_{\rm out}$ qualifying leaves. Suppose that its edges can be navigated
and its leaves tested reversibly in $2^{o(n)}$ work per operation.
Then all qualifying leaves can be reported with bounded error in
\[
 2^{o(n)}\sqrt{S(D_{\rm out}+1)}
\]
work, apart from separately charged index construction and output writes.
One qualifying leaf can be found, or their absence established, in
$2^{o(n)}\sqrt S$ work. The tree size need not be known beforehand.
Auxiliary workspace is $2^{o(n)}$, apart from the stored access data and
reported output records.
\end{lemma}
The tree must include the vertices needed to reach candidates, including
queries whose buckets are empty; counting only accepted pairs is
insufficient. For random trees the pointwise bound and Cauchy--Schwarz give
\[
 \mathbb E\sqrt{S(D_{\rm out}+1)}
 \leq\sqrt{\mathbb ES\;\mathbb E(D_{\rm out}+1)}.
\]
No independence between these two counts is needed. This lemma gives a
reporting algorithm with measured output records. It does not give a
canonical coherent superposition of all output lists, which is an
additional requirement for the QRAQM setup.

\begin{lemma}[Product-Johnson spectral gap
  {\cite[Definition~4 and Theorem~2]{BBSS}}]
\label{lem:johnson-gap}
Let $J(N,R)$ have the $R$-subsets of $\{1,\ldots,N\}$ as vertices;
one step exchanges a uniformly chosen member with a uniformly chosen
nonmember. Fix a constant number $q$ of factors. On the Cartesian product
$\prod_{i=1}^qJ(N_i,R)$, choose one factor uniformly for the exchange and
make the walk lazy with probability $1/2$.
If $R=o(\min_iN_i)$, its spectral gap is $\Theta(1/R)$.
In particular, $R=2^{hn+o(n)}$ gives gap $2^{-hn+o(n)}$.
\end{lemma}
For one factor the exchange-chain gap is $N/[R(N-R)]$; taking a constant
number of factors and adding laziness changes it by constant factors.
Deterministic auxiliary data attached canonically to every product vertex
do not change this underlying Markov chain. Whether such data can be
maintained coherently at the claimed cost is a separate algorithmic
requirement.

\begin{lemma}[Quantum-walk search
  {\cite[Theorem~3]{MNRS}}]
\label{lem:mnrs}
Let $P$ be a finite reversible ergodic Markov chain with spectral gap
$\delta>0$. Suppose that, whenever a marked state exists, its stationary
marked mass is at least $\varepsilon>0$.
Assume actual coherent implementations of stationary setup with canonical
state data, the forward and reverse walk updates and their inverses, and
the marked-state phase test. If these have uniform costs
$T_{\rm setup},T_{\rm up},T_{\rm check}$, respectively, then a marked
state can be found, or their absence detected, with bounded error at cost
\[
 O\!\left(T_{\rm setup}+
 \frac{1}{\sqrt\varepsilon}
 \left(\frac{T_{\rm up}}{\sqrt\delta}+T_{\rm check}\right)\right).
\]
\end{lemma}
This is the fixed-cost walk theorem. In particular it does not replace
$T_{\rm up}$ by the expected cost of an ordinary classical update, and it
does not by itself give a hard workspace bound. Our QRAQM update assumption
asks for an operational effective-update analogue of this scaling and
states its workspace requirement explicitly. A negligible checking rate
also requires that the canonical marked flag be maintained by the charged
updates. The displayed theorem explains the gap and marked-mass factors;
the stronger update interface remains conditional.



\section{Shared representation framework}\label{sec:framework}
\subsection{Compatible types and modular constraints}
In a binary representation tree, every parent vector is the integer
coordinatewise sum of its children. The root type assigns mass $1/2$
to each of $0$ and $1$. Child digits may be negative or repeated.
A useful law $Q(X,Y)$ prescribes the full joint type of two children;
its marginals and sum law must equal their child and parent types.
The useful predicate tests $Q$, not just the parent sum type.

Partition the coordinates into classes $C_w$ of fixed rational proportions
$\lambda_w$, and prescribe a type within each class. Write $W$ for the
class of a uniform coordinate. All logarithms are base two and
$0\log0=0$. By Lemma~\ref{lem:type-size}, a fixed-alphabet type class has size
\[
 |T_n(p)|=\frac{n!}{\prod_a(np(a))!}
       =2^{n\HH(p)+O(\log n)}.
\]
The classwise product has entropy rate $\HH(X\mid W)$.
Thus every entropy or information quantity below is conditioned on $W$;
we average its classwise values before taking a cost maximum.

\begin{lemma}[Useful-pair counts]\label{lem:type-counts}
The fraction of independent uniform marginal-type vectors with joint
law $Q$ is $2^{fn+O(\log n)}$, where
\[
 f=\HH(X,Y\mid W)-\HH(X\mid W)-\HH(Y\mid W)
   =-I(X;Y\mid W).
\]
For a fixed parent $s$ of the required sum type, the number of its
useful decompositions has rate
$\HH(X,Y\mid W)-\HH(X+Y\mid W)$.
\end{lemma}
\begin{proof}
The number of ordered pairs of type $Q$ is the multinomial coefficient
$n!/\prod_{a,b}(nQ(a,b))!$. Divide by the two marginal type-class sizes.
For a fixed parent, distribute the pair symbols among its positions of
each sum symbol $u$, giving
$\prod_u (nQ_{X+Y}(u))!/\prod_{a+b=u}(nQ(a,b))!$.
Apply Stirling's formula separately in every coordinate class.
\end{proof}
When a symmetric law is encoded as unordered masses $q_{\{a,b\}}$,
its ordered masses are $q_{\{a,b\}}/2$ off the diagonal. In particular
$\HH(Q)=\HH(q)+\sum_{a<b}q_{\{a,b\}}$; this orientation term is retained
in the entropy calculations.

A cumulative modular rate $c$ means a constraint modulo
$2^{cn+o(n)}$ on the scalar product of the whole vector with $a$.
Offsets are chosen from the root downwards. At each split the left
residue is uniform at the children's common modulus, and the right
residue is the parent's residue minus it at that modulus. Distinct
splits use independent seeds. These offsets and the filter family are
fixed during an execution or quantum walk. Source-support partitions
are refined with the coordinate classes and carry the corresponding
proportional types. A random coordinate permutation makes the planted
binary vector balanced in each fixed cell with only a polynomial loss.

Exact probability tables are used on admissible lengths. The required
divisibility modulus is constant because the profiles are fixed. For
other lengths, a bounded number of coordinates can be guessed separately;
a random permutation makes the remaining root balanced with constant
probability. The resulting bounded difference in modulus bit length
costs only a constant factor. Rounding realized list sizes and modular
bit lengths contributes $o(n)$ to the rates, tolerated by the fixed
positive resource margins.

\subsection{Typed filters and their implementation}
A filter is an extension $R(X,Y,Z\mid W)$ of a useful law $Q$.
For a center of type $R_Z$, query and data incidence test the prescribed
$XZ$ and $YZ$ marginal types. Define
\[
 \alpha=I(X;Z\mid W),\quad \beta=I(Y;Z\mid W),\quad
 \kappa=I(XY;Z\mid W),
\]
\[
 d=I(X;Z\mid Y,W),\quad q=I(Y;Z\mid X,W),\quad
 r=I(X;Y\mid Z,W).
\]
The chain rule gives
\begin{equation}
 d=\kappa-\beta,\quad q=\kappa-\alpha,\quad
 \kappa-\alpha-\beta=f+r,
 \qquad \kappa=-f+d+q-r.
 \label{eq:filter-identities}
\end{equation}
A constant label has $\alpha=\beta=\kappa=d=q=0$, while $r=-f$.

\begin{lemma}[Typed filter cover and static access; product construction following {\cite[Proposition~5]{HG26}}]
\label{lem:filter-access}
For fixed finite rational $R$, subexponential verified preprocessing
constructs an implicitly addressed filter family with at most
$2^{\kappa n+o(n)}$ addresses. It covers every useful pair and has at most
$2^{qn+o(n)}$ query incidences and $2^{dn+o(n)}$ data incidences per vector.
Incidences can be enumerated and a pair's first common address tested
with $2^{o(n)}$ overhead. A static list of $N$ data occurrences has an
index of $N2^{dn+o(n)}$ records with exact reversible bucket counts and
selection by rank. For static child lists, let $N_X$ be the number of
query occurrences, $I_X$ their total incident addresses and $C$ the
actual bucket-candidate visits. The candidate domain has a fixed binary
tree of polynomial depth and size at most
$2^{o(n)}(1+N_X+I_X+C)$, with reversible edge navigation and leaf tests
costing $2^{o(n)}$. These claims concern static navigation, not the
occupancies of recursively generated lists or coherent dynamic updates.
\end{lemma}
\begin{proof}
We give the construction within one coordinate class; a fixed number
of classes is handled by products. Split it into blocks of length
$b=\lceil(\log_2 n)^2\rceil$, allowing one shorter block.
Write $B(z\mid x,y)=R(x,y,z)/Q(x,y)$ on positive cells. Enumerate the
integer block-table patterns $m_t(x,y,z)$ with global sum $nR$ and
\[
 |m_t(x,y,z)-m_t(x,y)B(z\mid x,y)|\le1.
\]
There are at most $(b+1)^{O(n/b)}=2^{o(n)}$ patterns.
Every global useful pair admits one: for each $(x,y)$, the fractional
matrix $m_t(x,y)B(z\mid x,y)$ has integral row and column margins.
After taking floors, an integral bipartite flow rounds its residual
entries within one while preserving those margins. Zero cells remain zero.
Thus the construction covers global types without imposing one fixed
block pattern on every useful pair.

For a pattern let $R_t=m_t/n_t$ and $\overline R_t=(R_t)_{XY}B$.
Their averaged entropy features differ by $o(1)$, uniformly over patterns.
Under $\overline R$, the block index $T$ satisfies $Z\perp T\mid XY$,
so conditioning gives
\[
 I(XY;Z\mid T)\le\kappa,\quad
 I(X;Z\mid Y,T)\le d,\quad I(Y;Z\mid X,T)\le q.
\]
These inequalities bound the sums of local family and incidence exponents.

For each local integer table, let $p_{XY}$ be the exact probability that
a uniform local center gives its triple type for a fixed useful pair;
let $p_X,p_Y$ be its marginal incidence probabilities. Sample
$\lceil b^3/p_{XY}\rceil$ center IDs, permitting repeated strings.
A fixed pair is missed with probability at most $e^{-b^3}$.
Because $p_X,p_Y\ge p_{XY}$, Chernoff bounds give uniform incidence
bounds of twice their expectations plus a constant. A union bound over
all local strings, pairs and polynomially many tables succeeds.
Exhaustive local enumeration verifies coverage and incidence bounds in
$2^{O(b)}=2^{o(n)}$ time and space; failed samples are discarded.
Exact conditional type counts show local family rate $n_t\kappa_t$
and query/data incidence rates $n_tq_t,n_td_t$, up to $O(\log b)$.

Take the product of local codebooks for each pattern and the disjoint
union over patterns. An address is its pattern and local ID tuple.
The preceding bounds prove the three global rates. Refined local
incidences imply the original global marginal predicates, and the
rounding argument ensures every useful pair has a common address.
Mixed-radix enumeration of local incidence lists gives the stated
navigation. To select the first common address, scan the subexponentially
many patterns and their local codebooks for their least common IDs.
This costs $2^{o(n)}$ and has a clean reversible implementation.

For a static data list, enumerate incidences and sort records by pattern,
center tuple, modular key and occurrence ID. Every nonempty bucket is
an interval; binary searches find its endpoints, subtraction gives its
count and an offset implements selection by rank. Empty buckets return
zero and a dummy value. All keys have polynomial bit length. Read-only
queries and their inverses require polynomially many QRACM reads; sorting
and storage cost $N2^{dn+o(n)}$. No aggregate table over query prefixes is
constructed.

The candidate tree first chooses a query occurrence, then a pattern and
its local center IDs, and finally a rank in the actual terminal bucket.
Encode these choices by binary rank trees. Each stage has polynomial
bit length, and local incidence lists or the terminal bucket endpoints
supply its children reversibly. Empty patterns and buckets are retained
as rejecting headers. The vertex count is at most a polynomial factor times
$1+N_X+I_X+C$, including rejecting and empty query-address headers. Testing a leaf
checks its useful type and first common address. Thus the resulting
fixed static tree satisfies Lemma~\ref{lem:tree-reporting}; no totals below
a partially specified query-center prefix are needed.
\end{proof}
The incidence predicates may be refined by these block patterns.
The list/load assumption below applies to this actual construction.
Its exponentially large global family charge $\kappa$ is conservatively
retained in every resource bound. Canonical centers remove duplicate
filter witnesses, but not different source pairs producing the same
parent vector. Quantum reporting, actual padded loads and coherent
canonical setup have the separate interfaces stated next.

\subsection{Algorithmic assumptions}
\label{sec:assumptions}
Exact type counts and the static filter construction do not make the
recursive lists independent. The following assumptions specify the
additional properties used in the three proofs. They concern only the
three fixed profiles and the source, offset and filter constructions
used here. They are not assertions about every admissible probability law.

\begin{assumption}[Recursive lists, loads and survival]
\label{ass:lists}
For the specified construction, conditioning a type domain of rate $h$
on $c$ modular bits per coordinate gives the modeled rate $h-c$.
A useful-pair predicate contributes its exact exponent $f$, and a fresh
modular condition of rate $\Delta$ contributes $-\Delta$ to the merge
count. These estimates hold for the actual recursively produced lists,
jointly with their filter incidences, with at most $2^{o(n)}$ loss.
All candidate counts include unsuccessful candidates. Claimed occupancy
bounds can be enforced with subexponential overhead without losing an
exponentially larger fraction of solutions than of ordinary states.

For classical and QRACM search, the prescribed represented root survives
with probability $2^{-o(n)}$ after the charged search. For QRAQM, the
marked fraction of the product-subset space has rate $E$ derived from
its root recurrence. Required source-pair/parent occurrence distortions
are subexponential. These statements hold over fresh seeds for each
input in a set of probability $1-o(1)$, so repetition on a fixed good
input is permitted.
\end{assumption}
The survival clause is stronger than a first-moment calculation.
The capacities checked below are necessary domain and prefix conditions;
they do not by themselves prove survival or concentration.

\begin{assumption}[Access, padding and output reporting]
\label{ass:access}
The actual indices and query trees support the operations used in the
cost derivations with $2^{o(n)}$ overhead. Their candidate and padded
bucket domains have the stated rates, including empty addresses and
all unsuccessful candidates. Occurrence and canonical-address tests,
output writes and deduplication have the costs charged below.
Classical enumeration visits the actual incidence and candidate records.

In the quantum models, controlled navigation and tests are reversible.
For fixed static lists, all-output reporting and one-output search use
Lemma~\ref{lem:tree-reporting}; index construction and output writes are
charged separately. The actual tree size and accepted occurrence counts
must have the rates supplied by the list/load assumption. The implicit
QRACM routine additionally has the uniform padded occurrence states and
success probabilities, or coherent estimates sufficient for the chosen
amplification, needed to invoke Lemma~\ref{lem:amplitude-amplification}
within the charged costs. Rank/select in complete buckets alone does
not establish these preparation and padding properties.

For QRAQM, source preparation, bottom matching and all-output setup
must have a canonical coherent implementation on superposed source
subsets, with the charged setup time and workspace. Reporting and
measuring classical outputs alone would not provide that setup state.
Quantum errors are reduced sufficiently for the complete algorithmic composition, with only $2^{o(n)}$ overhead;
auxiliary storage is included in the charged bounds.
\end{assumption}
Lemma~\ref{lem:filter-access} supplies the typed cover, incidence
navigation and static exact bucket access; Lemma~\ref{lem:tree-reporting}
supplies the measured reporting bound for the resulting fixed tree.
The classical indexing, sequential bucket scans and deduplication are supplied by the construction once their actual counts are given.
The claimed sizes and output multiplicities for actual recursive lists,
QRACM padded state preparations and QRAQM coherent canonical setup remain
the corresponding assumptions above.

\begin{assumption}[Coherent expected-update walk interface]
\label{ass:updates}
For the specified QRAQM profile, the maintained representation is a
canonical function of the current source subsets. An exchange supports
the clean controlled transition, inverse, root test and extraction
operations required by a product-Johnson walk. The typed update costs
used in Section~\ref{sec:qraqm} are valid as effective coherent costs
under the actual conditioned states. With source-subset rate $H$,
marked rate $E$ and charged effective-update rate $U$, detection and
extraction have rate $(H-E)/2+U$, in addition to setup.
Uncomputation, composition and error control add no exponential factor.
Transition, marking and extraction have a hard workspace rate $w$ on
every represented branch; the advertised bound requires $w\le C$.
\end{assumption}
The marked-mass and gap factors agree with Lemma~\ref{lem:mnrs} for fixed-cost coherent updates. The operational effective-update clause above is an additional quantum implementation hypothesis; that lemma does not establish it. An ordinary
mean number of changed entries does not prove the coherent walk
interface or its worst-branch workspace bound. The QRAQM theorem keeps
it explicit; classical and QRACM search do not require it.

\subsection{Finite profiles and verification}
The supplementary verification software consists of three self-contained
Python programs: \path{verify_classical.py}, \path{verify_qracm.py} and
\path{verify_qraqm.py}. Each includes the complete exact rational
probability laws and finite rate prescriptions for its model. The
programs reconstruct the parent/child types and label couplings, check
their exact normalization and marginal identities, and evaluate the
entropy formulas with $100$-digit directed interval arithmetic.

The resource checks include every source, recurrence, capacity, family,
incidence, candidate, reporting and final search term in the corresponding
analysis. The QRAQM program also checks the exact common-height
prescription, the information identities and the source-domain and
storage budgets. The three finite constructions are fixed inputs to
these checks; reproducing the bounds requires no numerical optimization.
The shorter decimal tables in the text give the scale of the resource
bounds, while the exact embedded parameters and interval comparisons
certify the strict inequalities.

The mathematical arguments establish the three resource bounds from
the stated algorithmic assumptions. The numerical checks establish the
probability identities and resource inequalities for the chosen profiles.
They do not establish the list/survival, access or coherent-update
assumptions. In particular, proving the complete higher-level coherent
update remains necessary for an unconditional QRAQM algorithm running in time and total space
$2^{0.198n+o(n)}$.



\section{Classical filtered representation trees}
\label{sec:classical}
The classical algorithm builds every list in a binary representation tree.
At each merge a type filter reduces the pairs that must be tested, while
modular constraints reduce the number of partial sums that are retained.
Both costs are charged: the algorithm writes its filter incidences and
visits every matching candidate, including candidates whose joint type
is not useful.

\subsection{Construction and complete merge cost}
The tree has depth $D=7$ and digit alphabet
$\mathcal A=\{-2,-1,0,1,2,3\}$. There are two equally sized coordinate
classes. At the root their ordered useful-pair laws are transposes of one
another. The two prescribed subtree profiles are exchanged between the
classes in the opposite root subtree. Within a subtree, child exchange
is coupled to a permutation of the five filter labels. Consequently every
node at the same depth has the same whole-vector entropy rates, and
$\alpha_j=\beta_j$. This statement concerns complete vectors: entropy
rates are averaged across the classes before any resource maximum is
taken. There is one shared modulus at each depth.

Let $h_j$ be the entropy of a child type at depth $j$, and let $Q_j$ and
$R_j$ denote the class-conditioned useful-pair law and its filter
extension. Write
\[
 f_j=\HH(Q_j\mid W)-2h_j,\qquad
 i_j=\kappa_j-\alpha_j=\kappa_j-\beta_j,\qquad
 r_j=I(X;Y\mid Z,W).
\]
The actual marginal and sum laws in the two classes agree at every edge;
entropy identities alone would not ensure this compatibility. The root
sum type is binary with equal masses, so $h_0=1$.

Put $c_0=1$, and let $c_j$ be the cumulative modular rate at depth $j$,
with $1\ge c_1\ge\cdots\ge c_D\ge0$. Sibling residues are chosen so that
their sum agrees with their parent's residue on the already imposed bits.
A merge from depth $j$ then checks the additional rate
$\Delta_j=c_{j-1}-c_j$. If each child list has rate $\ell_j$, the list/load
contract gives
\begin{equation}
 \ell_{j-1}=2\ell_j-\Delta_j+f_j,\qquad \ell_0=0.
 \label{classical:recurrence}
\end{equation}
The factor $2^{f_jn+o(n)}$ comes from exact joint-type counting; the modular
factor for recursively generated lists is part of
Assumption~\ref{ass:lists}.

To build a bottom list, divide each coordinate class into two equal
supports. Sample $2^{un+o(n)}$ vectors of the prescribed type on each
support, and match their scalar products modulo $2^{c_Dn+o(n)}$.
The source population on either support has entropy $h_D/2$, and
\begin{equation}
 \ell_D=2u-c_D,\qquad
 u\le h_D/2,\qquad
 \ell_j+c_j\le h_j\quad(1\le j\le D).
 \label{classical:capacity}
\end{equation}
Under the source-load contract, a residue dictionary constructs this list
in time and space $2^{\max\{u,\ell_D\}n+o(n)}$. Every bottom source and
list is included; the constant number of tree positions changes no
exponent.

For an internal merge, Lemma~\ref{lem:filter-access} provides the
class-conditioned cover and its incidence navigation. Conservatively
constructing the entire family has rate $\kappa_j$. The two input-incidence
tables each have rate $\ell_j+i_j$. Scan their modularly matching
addresses and test the useful-pair predicate, reporting a useful pair at
its first common filter address. The candidate-address rate is
\begin{align}
 2\ell_j-\Delta_j+\kappa_j-\alpha_j-\beta_j
 &=\ell_{j-1}-f_j+\kappa_j-\alpha_j-\beta_j\notag\\
 &=\ell_{j-1}+r_j.
 \label{classical:candidates}
\end{align}
The second equality is the information identity
$\kappa_j-\alpha_j-\beta_j=f_j+r_j$.
It includes candidates of the wrong useful type and repeated center
addresses. Canonical address tests remove repeated center witnesses;
coincident output vectors are deduplicated. Distinct-output sizes and
candidate loads for these correlated lists still use the stated list,
occurrence and access contracts.

Input scans are bounded by the incidence tables since $i_j\ge0$, and
output writes are bounded by the candidate term since $r_j\ge0$.
Candidates can be processed sequentially. Thus all sources, stored lists,
families, incidence writes and candidate tests are bounded by
\begin{equation}
 \tau_C=\max\left\{u,
       \max_{1\le j\le D}
          \{\kappa_j,\ell_j+i_j,\ell_{j-1}+r_j\}\right\}.
 \label{classical:cost}
\end{equation}
This is both a time bound and a space bound for the construction, under
Assumptions~\ref{ass:lists} and~\ref{ass:access}.

\subsection{A finite classical construction}
The exact useful and filter probabilities are embedded in
\path{verify_classical.py}. The two class proportions are $1/2$, and the
classwise sum and marginal identities hold over the rationals. The
class-$L$ level-six filter and both level-seven filters have a constant
label. In particular $i_7=0$ exactly. These laws and the following formulas
define one finite construction; no limiting parameter choice is needed.

Set
\begin{align}
 t_0&=\frac{1-h_4-\sum_{j=1}^{4}f_j+
                         \sum_{j=1}^{4}i_j}{4},
 &\delta&=10^{-6},
 &t&=t_0+\delta/4,
 \label{classical:time-profile}\\
 \ell_j&=t-i_j\quad(1\le j\le7),
 &c_1&=1-2\ell_1-f_1,
 &c_j&=c_{j-1}+\ell_{j-1}-2\ell_j-f_j\quad(j\ge2),
 \label{classical:rates-profile}\\
 u&=(\ell_7+c_7)/2.
 \label{classical:source-profile}
\end{align}
The interpretation is simple: every input-incidence table has the same
rate $t$. The value of $t_0$ is chosen so that their common cost balances
the capacity at depth four. Indeed summing the modulus recurrence gives
\[
 c_4=1-\sum_{j=1}^{3}\ell_j-2\ell_4-\sum_{j=1}^{4}f_j,
 \qquad h_4-c_4-\ell_4=\delta.
\]
The same recurrence makes every merge count in
\eqref{classical:recurrence} an exact identity, including the root count.
There are no separately tuned list heights or modulus rates.

\begin{center}\small
\begin{tabular}{@{}crrr@{}}\toprule
Depth $j$ & List rate $\ell_j$ & Modulus rate $c_j$ & Candidate rate\\\midrule
1 & 0.211967423269 & 0.999982060603 & 0.263380572935\\
2 & 0.186940098937 & 0.999980726081 & 0.263382504052\\
3 & 0.196177821947 & 0.999976521029 & 0.263383204185\\
4 & 0.203067064704 & 0.956415994747 & 0.263383440534\\
5 & 0.236032087788 & 0.778104141382 & 0.263372923914\\
6 & 0.258822538842 & 0.526748054134 & 0.263377914550\\
7 & 0.263383502113 & 0.263351278992 & 0.263370229084\\\bottomrule
\end{tabular}
\end{center}
The table is display rounding only. Evaluation of the exact laws gives
\begin{align*}
 t_0&=0.2633832521129001557945171803\ldots,\\
 t&=0.2633835021129001557945171803\ldots,\qquad
 u=0.263367390552\ldots .
\end{align*}
Directed-interval verification checks nonnegative list rates, all modulus
bounds and nesting, the leaf-domain bound, every capacity in
\eqref{classical:capacity}, and every family and candidate term in
\eqref{classical:cost}. The capacities and modulus gaps have fixed positive
reserves; their minimum is at least $10^{-6}$. Every incidence rate is
exactly $t$, and all remaining charged costs are at most $t$. The strict
gap to $0.26339$ exceeds $6.4978\cdot10^{-6}$. The exact probability tables,
coefficient identities and outward-rounded interval comparisons are
checked by the verification program; printed decimal truncation is not
used as a certificate.

\begin{proof}[Proof of Theorem~\ref{thm:classical}]
The exact type-flow checks give a compatible whole-vector representation
tree, and \eqref{classical:rates-profile} gives valid nested modular
constraints and the list recurrences. The capacities ensure that no
modeled list exceeds its underlying type domain. The source construction
and the merge analysis leading to \eqref{classical:cost} account for all
work and persistent storage; their maximum is at most
$t<0.26339$.

To identify the survival calculation used by the list contract, put
$s_j=\ell_j+c_j-h_j$ and
$\rho_j=\HH(Q_j\mid W)-h_{j-1}$. Type compatibility and the recurrences
imply
\[
 s_D=2u-h_D,\qquad s_{j-1}=\rho_j-c_j+2s_j.
\]
All $s_j$ are nonpositive by the checked capacities, and $s_0=0$.
Thus the corresponding capped reference recursion has root inclusion
rate zero. Assumption~\ref{ass:lists} is used here to transfer that
survival calculation to the actual correlated recursive lists, giving
success $2^{-o(n)}$ on its good-input set. Fresh-seed repetition yields
constant success without changing the exponent. Every reported root is
checked to be binary and to satisfy the target equation. The entropy
identities and capacity checks establish the stated accounting; they do
not themselves prove the survival or modular-load assumption.
\end{proof}



\section{Quantum search with classical memory}
\label{sec:qracm}

The QRACM algorithm stores one side of each search classically and
generates the other side coherently.  Only a single root-to-leaf path is
implicit.  Its sibling subtrees are constructed once, after which nested
amplitude amplification searches the path.  Every stored merge must
report all its outputs; only the implicit path uses state preparation.

\subsection{The tree and its sources}
Use a binary tree of fixed depth $D$, indexed from root $0$, with children
$v_L=2v+1$ and $v_R=2v+2$.  Write
\[
 \mathcal P=\{0,1,3,\ldots,2^D-1\},\qquad b_\star=2^D-1.
\]
All positive-depth nodes outside $\mathcal P$ are stored.  The classwise
ordered law $Q_v(X,Y\mid W)$ specifies the useful pairs at an internal
node: its marginals are the child types and its sum law is the parent
type.  Let
\[
 h_v=H(P_v\mid W),\qquad
 f_v=H(Q_v\mid W)-h_{v_L}-h_{v_R}.
\]
A list at $v$ has rate $\ell_v\ge0$ and cumulative modular rate
$0\le c_v\le1$.
Siblings have equal moduli.  With $c_0=1$, $\ell_0=0$ and
$\Delta_v=c_v-c_{v_L}\ge0$, the required rates satisfy
\begin{align}
 \ell_v&=\ell_{v_L}+\ell_{v_R}-\Delta_v+f_v,
 \label{qracm:recurrence}\\
 \ell_v+c_v&\le h_v\quad(v>0),&
 \ell_{b_\star}+c_{b_\star}&=h_{b_\star}.
 \label{qracm:capacity}
\end{align}
The recurrence combines the exact type-counting exponent $f_v$ with the
modular-load hypothesis.  The capacities are imposed at implicit nodes
as well as stored nodes: a full bottom source alone does not guarantee
that a represented solution has a sufficiently likely earlier prefix.

For every stored bottom list, split each coordinate class into equal
complementary supports.  Sample and store a uniform subset of rate
$(\ell_v+c_v)/2$ from each half-type domain, using ordinary type ranking
and classical dictionaries.  The capacity in \eqref{qracm:capacity}
ensures that the requested subsets fit their domains.  Index the two
tables by residue and report all matching pairs; their cost is the
maximum of the source and output rates under the modular-load condition.
For
the implicit bottom source, use support fraction $\rho=1/3$ for the
stored right domain, of rate
\[
 b=\rho h_{b_\star}.
\]
Prepare the complementary source uniformly and query this right
dictionary.  The support cells carry their prescribed marginal types;
they are not unrestricted global type classes.  Under the actual bucket
padding and modular-load conditions, Lemma~\ref{lem:amplitude-amplification} applied to the
bottom modular equation prepares a uniform occurrence state in time
$2^{g_{b_\star}n+o(n)}$, where
\begin{equation}
 g_{b_\star}=\tfrac12[c_{b_\star}-b]_+.
 \label{qracm:bottom-generator}
\end{equation}
Indeed the padded right bucket has rate $[b-c_{b_\star}]_+$, and the
marked fraction has rate $-[c_{b_\star}-b]_+$.  The construction and
storage of the full right source, of rate $b$, are also charged.

\subsection{Stored merges and implicit amplification}
At node $v$ let $\alpha_v,\beta_v,\kappa_v$ be the filter information
rates from Lemma~\ref{lem:filter-access}.  Query addresses use $XZ$;
the stored data dictionary uses $YZ$.  Its incidence rate is
$\ell_{v_R}+\kappa_v-\beta_v$.  The complete candidate tree for a stored
merge has rate
\begin{equation}
 C_v=\max\{\ell_{v_L}+\kappa_v-\alpha_v,
 \ell_{v_L}+\ell_{v_R}-\Delta_v+
               \kappa_v-\alpha_v-\beta_v\}.
 \label{qracm:candidates}
\end{equation}
The first term counts query addresses, including empty buckets; the
second counts every modular candidate, including wrong useful types.
Mark a useful pair only at its first common filter address.  The
occurrence hypothesis bounds the number of distinct pairs representing
one output vector by $2^{o(n)}$.  Consequently the number to report has
rate $\ell_v$. Lemma~\ref{lem:tree-reporting}, applied to the fixed static candidate tree, therefore gives all-output reporting cost
$2^{(C_v+\ell_v)n/2+o(n)}$.  The output vectors are measured and inserted
into a classical dictionary.  Input scans, output writes, data-incidence
writes and the covering-family charge give the complete stored-merge
rate
\begin{equation}
 s_v=\max\{\ell_{v_L},\ell_{v_R},\ell_v,\kappa_v,
       \ell_{v_R}+\kappa_v-\beta_v,(C_v+\ell_v)/2\}.
 \label{qracm:stored-cost}
\end{equation}
The output term is essential: charging only the search for one useful
pair would not construct the stored subtree.

Along the implicit path, assume the left child is already prepared
uniformly over its occurrence addresses, with cost exponent $g_{v_L}$.
For each incoming occurrence use the same padded rectangle of query
addresses and right-bucket slots, of rate
\[
 \chi_v=\max\{\kappa_v-\alpha_v,
       \ell_{v_R}-\Delta_v+\kappa_v-\alpha_v-\beta_v\}.
\]
The useful partner count per incoming occurrence has modeled average
rate $\lambda_v=\ell_{v_R}-\Delta_v+f_v$.  Thus the accepted fraction
has rate $\lambda_v-\chi_v$, and Lemma~\ref{lem:amplitude-amplification}, under the stated preparation interface, gives
\begin{equation}
 g_v=g_{v_L}+\tfrac12\max\{
   \kappa_v-\alpha_v-\ell_{v_R}+\Delta_v-f_v,
   \kappa_v-\alpha_v-\beta_v-f_v\}.
 \label{qracm:generator}
\end{equation}
The increment is nonnegative, since its second term is
$I(X;Y\mid Z,W)\ge0$.  Canonical filter addresses remove repeated
filter witnesses.  Other source-pair occurrences remain in the state;
the occurrence and survival conditions in Assumptions~\ref{ass:lists}
and~\ref{ass:access} apply to these actual states.  They cannot be
replaced by a calculation for independent random list entries.

The algorithm now has four steps: construct the stored bottom lists;
construct their stored ancestors using \eqref{qracm:stored-cost}; build
the right data indices along $\mathcal P$; and apply
\eqref{qracm:bottom-generator} and \eqref{qracm:generator} from the bottom
to the root.  Measure a root occurrence and verify its binary coordinates
and target sum.  The list/survival condition justifies the prescribed
subexponential repetitions over public seeds.

Let $\mathcal S$ be the stored nodes and $\mathcal B$ the stored bottom
nodes; a superscript $\circ$ restricts a set to its internal nodes.
The complete static construction and storage rate is bounded by
\begin{align}
 \sigma=\max\biggl\{&b,\ \max_{v\in\mathcal S}\ell_v,
          \max_{v\in\mathcal B}(\ell_v+c_v)/2,\notag\\
       &\max_{v\in\mathcal S^\circ}s_v,
          \max_{v\in\mathcal P^\circ}
          \{\kappa_v,\ell_{v_R}+\kappa_v-\beta_v\}\biggr\}.
 \label{qracm:setup}
\end{align}
The root filter is constant, so including its last two terms adds only
an already stored raw list.  Reporting workspace is included in its
charged setup cost.  Implicit state preparation retains a constant-depth
stack of addresses and vectors, rather than an exponential quantum list;
the reversible-access condition bounds its auxiliary overhead by
$2^{o(n)}$.  Hence the time and total-space rates are at most
$\max\{\sigma,g_0\}$ and $\sigma$, respectively.

\subsection{The finite profile and the constant}
\begin{proof}[Proof of Theorem~\ref{thm:qracm}]
Use the rational probability laws and rates embedded in
\path{verify_qracm.py}.  They have $D=4$, alphabet
$\{-1,0,1,2\}$, three filter labels and two equal coordinate classes.
The implicit path is $\{0,1,3,7,15\}$.  All fifteen stored bottom lists
have the exactly common rational rate
\[
 \ell_v=1123259302947/5349021095909
       =0.20999343296777855938\ldots .
\]
The root filter and filters at nodes $7,\ldots,14$ are constant.
The program includes every remaining rational child modulus.  The full-source identity and
root recurrence determine the missing implicit bottom rate and root
child modulus; the other list rates follow bottom-up from
\eqref{qracm:recurrence}.  Exact rational checks verify normalization,
every classwise marginal and sum law, and every label law.

Directed interval evaluation checks \eqref{qracm:capacity}, modulus
nesting, nonnegative rates and every term of
\eqref{qracm:stored-cost}--\eqref{qracm:setup}, including all eleven
stored internal merges and fifteen stored source constructions.
The principal resource bounds, rounded upwards for display, are
\begin{center}\small
\begin{tabular}{@{}lr@{}}\toprule
Resource & Rate upper bound\\\midrule
Implicit right-source construction & $0.169265570008$\\
Largest stored bottom-source construction & $0.209993773632$\\
Largest stored raw list & $0.209993432968$\\
Largest data-incidence index & $0.209994407551$\\
Largest all-output report & $0.209994407551$\\
Complete implicit root preparation & $0.209994407550$\\\bottomrule
\end{tabular}
\end{center}
The complete, unrounded maxima satisfy
\[
 \sigma,\max\{\sigma,g_0\}<0.209994407550290<0.2099947,
\]
with minimum strict target margin greater than
$2.9244971095\cdot10^{-7}$.  The program reconstructs the full resource expressions directly from
the ordered rational laws.  Applying the construction and amplification analysis
above proves the stated conditional bounds.  These arithmetic checks
certify the chosen profile; its recursive-list, actual-padding,
occurrence and survival conditions are precisely the stated assumptions.
\end{proof}



\section{A quantum-memory walk with one list height}
\label{sec:qraqm}

The QRAQM algorithm maintains a representation tree while a quantum walk
changes its sampled sources. Its useful simplification is that every
positive-depth list has the same size exponent. Its directed filters also
distinguish writes from searches: an entry is stored under its data
incidences, while its query incidences are generated coherently when
needed. This distinction gives the update costs below. The complete
coherent update, including propagation through dependent lists, remains
the additional premise in Assumption~\ref{ass:updates}.

\subsection{Sources and maintained lists}

Use a binary tree of depth $D$. There are two equally sized coordinate
classes. An ordered root law $Q_1$ has a balanced binary sum marginal;
its transpose is used in the other class. Below the root, two prototype
profiles are exchanged between the classes in the two root subtrees.
Within either prototype, the useful law $Q_j(X,Y)$ has equal child
marginals $P_j$ and sum marginal $P_{j-1}$. Its filter law extends this
useful law to $R_j(X,Y,Z)$. Thus all marginal and sum identities hold
separately in each class. All entries are complete vectors; there is no
separate algorithm on either class. Write
\[
 h_j=\HH(P_j\mid W),\qquad
 f_j=\HH(X,Y\mid W)-\HH(X\mid W)-\HH(Y\mid W).
\]
Here $f_1$ uses the ordered root law. For $j\ge2$, the displayed
quantities average the two prototype laws with equal weights.

Let $c_0=1$ be the target-modulus rate and let $c_j$ be the cumulative
modular rate at depth $j$. Put $\ell_j=H$ at every positive depth and
write $E=\ell_0$ for the root-completion exponent. The list model is
\begin{equation}
 E=2H-1+c_1+f_1,\qquad
 H=2H-c_{j-1}+c_j+f_j\quad(2\le j\le D).
 \label{qraqm:recurrence}
\end{equation}
Type counting gives the useful-pair factor $2^{f_jn+o(n)}$; the modular
factor in this recurrence for dependent lists uses
Assumption~\ref{ass:lists}. The finite profile must satisfy
\begin{equation}
 1\ge c_1\ge\cdots\ge c_D\ge0,\quad H>0,\quad E\le0,
 \qquad H+c_j-E/2^j\le h_j\quad(1\le j\le D).
 \label{qraqm:capacities}
\end{equation}
The last inequalities are the reverse-prefix capacities; they do not
by themselves establish survival of a prescribed root representation.

At each bottom position, split the support in proportion
$\rho=c_D/h_D$ within each coordinate class. Enumerate the complete
right-source type domain, of rate $c_D$, and index it by its modular
residue. Its complementary left domain has rate
$\eta=h_D-c_D$. A walk vertex chooses an independent uniform subset
of $m=2^{Hn+o(n)}$ left restrictions at each bottom position. Matching
the selected restrictions with the right source gives the modeled
bottom lists. The profile has $\eta>H$, as required for these subsets.
The fixed right tables cost $2^{c_Dn+o(n)}$; explicitly allowing the
full left domain as a carrier costs $2^{\eta n+o(n)}$. Both charges are
included below. Source preparation, matching and canonical access use
Assumption~\ref{ass:access}; the complete source exchanges and propagated
changes use Assumption~\ref{ass:updates}.

\subsection{Setup and update costs}

At a filtered merge, let $X$ be a query entry and $Y$ a data entry. Put
\[
 \alpha_j=I(X;Z\mid W),\quad \beta_j=I(Y;Z\mid W),\quad
 \kappa_j=I(XY;Z\mid W),
\]
and use the equivalent information quantities
\[
 d_j=I(X;Z\mid Y,W),\qquad q_j=I(Y;Z\mid X,W),\qquad
 r_j=I(X;Y\mid Z,W).
\]
The chain rules give
\begin{equation}
 d_j=\kappa_j-\beta_j,\qquad
 q_j=\kappa_j-\alpha_j,\qquad
 \kappa_j-\alpha_j-\beta_j=r_j+f_j.
 \label{qraqm:information}
\end{equation}
Lemma~\ref{lem:filter-access} supplies a data-incidence index of rate
$H+d_j$ and $2^{q_jn+o(n)}$ query addresses per entry. Both sibling
lists have data indices; neither stores an additional query-side index.
When either sibling changes, its changed entry can play the query role
because the useful law is symmetric within each class. The filter must
cover both ordered orientations. A pair key is always tied to the two
tree positions, so reversed queries update the same parent occurrence.

The fresh modular rate is
$\Delta_j=c_{j-1}-c_j=H+f_j$. Hence the complete setup candidate tree has
rate
\[
 \chi_j=\max\{H+\kappa_j-\alpha_j,
       2H-\Delta_j+\kappa_j-\alpha_j-\beta_j\}
       =H+\max\{q_j,r_j\}.
\]
The first term charges empty query addresses, and the second charges
all bucket candidates, including candidates rejected by the useful-type
test. For fixed static lists, Lemma~\ref{lem:tree-reporting} reports all $2^{Hn+o(n)}$ parent occurrences at cost
$2^{(\chi_j+H)n/2+o(n)}$. The walk requires this construction coherently
with canonical outputs under Assumption~\ref{ass:access}; ordinary
quantum enumeration followed by measurements would not supply its setup
state. Including the family construction,
incidence writes and output storage, the setup rate at this merge is
therefore
\begin{equation}
 \max\{\kappa_j,H+d_j,H,H+q_j/2,H+r_j/2\}.
 \label{qraqm:merge-setup}
\end{equation}
The family charge $\kappa_j$ is retained for both time and storage.
Canonical first-common-address tests remove repeated filter witnesses.
They do not remove multiple source pairs for one parent vector; that
multiplicity is covered by the occurrence part of
Assumption~\ref{ass:lists}. Parent representation counts must be maintained
so that deleting one pair does not delete a still-represented vector.

For one changed entry, writing its data incidences has rate $d_j$ and
searching its candidate tree has rate $\max\{0,q_j,r_j\}/2$. This motivates
the internal update bound
\[
 B=\max_{2\le j\le D}\{0,d_j,q_j/2,r_j/2\}.
\]
There is no output-growth exponent in this expression because parent
and child heights are equal. Its use for a complete coherent update
requires Assumption~\ref{ass:updates}; an average degree at an isolated
merge is insufficient. The root is unfiltered. A reverse root query tests the transpose of the
ordered root law; averaging over coordinate classes does not make that
predicate symmetric. Its one-solution setup cost is the maximum of $H/2$ and $(2H-1+c_1)/2$, and its update
candidate rate is $(H-1+c_1)/2$, with a zero floor.

Writing $K=\max_{2\le j\le D}\kappa_j$, the complete setup and persistent
storage rates are consequently bounded by
\begin{align}
 \sigma&=\max\{c_D,\eta,K,H+B,(E-f_1)/2\},
       \label{qraqm:setup}\\
 \mu&=\max\{c_D,\eta,H,K,\max_{j\ge2}(H+d_j)\}\le\sigma.
       \label{qraqm:storage}
\end{align}
Setup workspace is charged to $\sigma$ by the access contract. A valid
declared update rate must satisfy
\begin{equation}
 U\ge B,\qquad U\ge(H-1+c_1)/2.
 \label{qraqm:update}
\end{equation}
At the root, a reverse query tests the transposed ordered root law;
averaging the coordinate classes does not make that predicate symmetric.

\subsection{The finite profile and the walk bound}

The exact rational laws embedded in \path{verify_qraqm.py} use $D=9$,
alphabet
$\{-2,-1,0,1,2\}$ and three filter labels. Filters at levels six through
nine are constant; their ordinary useful-pair constraints remain in place.
Define
\begin{align*}
 L_0&=-f_2=0.16360358795774631025\ldots,\\
 U_0&=d_2=0.03439616742177458384\ldots,\\
 T_0&=\kappa_3=0.19799975538308076878\ldots,
 \qquad \delta=2^{-26},
\end{align*}
and set
\begin{equation}
 \begin{split}
 H&=L_0+\delta,\qquad E=-L_0-2\delta,\qquad U=U_0+\delta,\\
 c_j&=1-(j+2)L_0-\sum_{i=1}^{j}f_i-(j+3)\delta,\qquad
 C=T_0+3\delta.
 \end{split}
 \label{qraqm:finite-rates}
\end{equation}
These formulas impose \eqref{qraqm:recurrence} exactly. In particular
$c_1-c_2=\delta$, so even this otherwise tight modulus increment is
strictly positive. No list height or modulus is an independently chosen
resource parameter.

The exact type flows, label laws and coefficient identities are checked
over the rationals. Directed interval evaluation of their entropy
expressions verifies every inequality in
\eqref{qraqm:capacities} and~\eqref{qraqm:update}, together with
$C\ge\sigma$ and $C\ge(H-E)/2+U$. The minimum nonstructural resource
slack exceeds $7.45414\cdot10^{-9}$. The following conservative rounded
bounds show the resulting scale; the exact formulas define the point.
\begin{center}\small
\begin{tabular}{@{}lr@{}}\toprule
Quantity & Upper bound\\\midrule
Fixed right source $c_D$ & $0.197999576479$\\
Full left carrier $\eta$ & $0.163923305055$\\
Setup and persistent storage $\sigma,\mu$ & $0.197999770281$\\
Declared walk budget $(H-E)/2+U$ & $0.197999792633$\\
Total declared rate $C$ & $0.197999800087$\\\bottomrule
\end{tabular}
\end{center}

\begin{proof}[Proof of Theorem~\ref{thm:qraqm}]
Construct the fixed source tables and the canonical representation tree,
including the data-incidence indices and all parent lists. The preceding
cost derivation gives setup exponent $\sigma$ and persistent storage
exponent $\mu$. All sources and merges are counted; the tree has a fixed
number of positions.

The graph is the product of the Johnson graphs on the sampled left
restrictions. Since $\eta>H>0$, Lemma~\ref{lem:johnson-gap} gives inverse spectral-gap exponent $H$.
Augmenting a vertex by a canonical function of its source sets does not
change this graph. The list/survival contract gives marked mass
$2^{En+o(n)}$ for the prescribed root completion.

Apply the operational coherent-update interface of
Assumption~\ref{ass:updates}, with effective update rate $U$. Its detection
and extraction guarantee has the inverse-square-root marked-mass and
spectral-gap scaling, including coherent controls and sufficiently small
composed error. The resulting walk exponent is $(H-E)/2+U$.
This step uses that full interface, not just a classical mean update
cost. Thus time is bounded by
$2^{\max\{\sigma,(H-E)/2+U\}n+o(n)}$. Setup workspace, persistent
storage and hard walk workspace give total space at most
$2^{\max\{\sigma,\mu,w\}n+o(n)}$. The finite interval checks bound both
maxima by $C<0.198$. The stipulated good-input survival and permitted
fresh-seed repetitions give constant success without changing an exponent.
The output is extracted and checked against the binary target equation.
Any $2^{o(n)}$ type-balancing or fresh-seed repetition allowed by the
list/survival contract is absorbed in the displayed bounds.
\end{proof}


\bibliographystyle{plainnat}
\bibliography{references}
\appendix

\section{Use of artificial intelligence}
ChatGPT was used extensively in developing this work, including
mathematical exploration, numerical calculations, and preparation and
revision of the manuscript. The authors take responsibility for the
mathematical claims, computations, references and final text.


\end{document}
